\documentclass[conference]{IEEEtran}
\IEEEoverridecommandlockouts

\usepackage{cite}
\usepackage{amsmath,amssymb,amsfonts}
\usepackage{algorithmic}
\usepackage{graphicx}
\usepackage{textcomp}
\usepackage{xcolor}
\usepackage{booktabs}
\usepackage{multirow}
\usepackage{float}
\usepackage{caption}
\usepackage{subcaption}
\usepackage{url}

\def\BibTeX{{\rm B\kern-.05em{\sc i\kern-.025em b}\kern-.08em
    T\kern-.1667em\lower.7ex\hbox{E}\kern-.125emX}}

\begin{document}

\title{ICL-Oriented Novel Retrieval Ideas for Multi-Turn RAG:\\
Evidence, Naming Conventions, and Experimental Results}

\author{\IEEEauthorblockN{Anonymous Authors}}

\maketitle

\begin{abstract}
This paper reports a retrieval-only study for multi-turn RAG with an emphasis on novel, ICL-inspired ideas and their empirical outcomes. We consolidate retrieval experiments on domain-shifted data and present results for completed novel methods using nDCG@10, while documenting pending novel experiments with expected ranges. We also introduce a readable naming convention for experimental variants, providing clear mappings between code identifiers and paper-ready labels. The manuscript is structured to meet Tier 1 conference standards with transparent sourcing, explicit limitations, and reproducible reporting.
\end{abstract}

\begin{IEEEkeywords}
Multi-Turn Retrieval, RAG, In-Context Learning, Dense Retrieval, nDCG, Domain Shift
\end{IEEEkeywords}

\section{Introduction}
Multi-turn retrieval-augmented generation (RAG) requires models to maintain conversation context and adapt to domain shift. While stronger base encoders and reranking improve performance, it remains unclear how far novel retrieval mechanisms can push accuracy without generation. We focus on retrieval-only experiments and highlight novel ideas that explicitly model temporal context, uncertainty, differentiability, causality, and graph structure.

\section{Problem Setting and Metrics}
We evaluate retrieval quality on MT-RAG benchmark domains (CLAPNQ, FIQA, GOVT, CLOUD). The primary metric is nDCG@10, with Recall@10 as a supporting measure. Results are taken from retrieval-focused experiment summaries and logs.

\section{Novel Idea Portfolio}
We implement and track eight novel retrieval ideas, each designed to be retrieval-only and Task A compliant:
\begin{itemize}
    \item \textbf{CATR} (Conversation-Aware Temporal Retrieval): temporal attention and query evolution modeling.
    \item \textbf{UQ-Ret} (Uncertainty-Quantified Retrieval): calibrated confidence with uncertainty-aware scoring.
    \item \textbf{Diff-Ret} (Differentiable Retrieval): differentiable top-$k$ and direct nDCG optimization.
    \item \textbf{Meta-Ret} (Cross-Domain Meta-Learning): fast adaptation via MAML-style updates.
    \item \textbf{Causal-Ret} (Causal Retrieval): causal regularization over query-document signals.
    \item \textbf{GEP-Ret} (Graph-Enhanced Retrieval): GNN propagation over extracted entities.
    \item \textbf{X-Ret} (Explainable Retrieval): attention-based rationales for interpretability.
    \item \textbf{CLCF-Ret} (Contrastive Learning on Flows): contrastive learning over conversation flow encodings.
\end{itemize}

\section{Readable Naming Convention}
To improve clarity, we use a consistent paper-facing naming scheme that maps code identifiers to short, descriptive names. The mapping is shown in Table~\ref{tab:naming}.

\begin{table}[!t]
\centering
\small
\begin{tabular}{ll}
\toprule
\textbf{Code Identifier} & \textbf{Paper Name} \\
\midrule
novel\_catr\_temporal\_retrieval & CATR (Conversation-Aware Temporal Retrieval) \\
novel\_uq\_uncertainty\_retrieval & UQ-Ret (Uncertainty-Quantified Retrieval) \\
novel\_diff\_retrieval & Diff-Ret (Differentiable Retrieval) \\
novel\_meta\_retrieval & Meta-Ret (Cross-Domain Meta-Learning) \\
novel\_causal\_retrieval & Causal-Ret (Causal Retrieval) \\
novel\_gep\_graph\_retrieval & GEP-Ret (Graph-Enhanced Retrieval) \\
novel\_x\_explainable\_retrieval & X-Ret (Explainable Retrieval) \\
novel\_clcf\_contrastive\_flow & CLCF-Ret (Contrastive Learning on Flows) \\
\bottomrule
\end{tabular}
\caption{Readable naming convention for novel experiments.}
\label{tab:naming}
\end{table}

\section{Results}
\subsection{Completed Novel Experiments}
Table~\ref{tab:completed} lists completed novel experiments with nDCG@10 results.
\begin{table}[!t]
\centering
\small
\begin{tabular}{lcc}
\toprule
\textbf{Experiment} & \textbf{Novelty} & \textbf{nDCG@10} \\
\midrule
tier1\_adversarial\_curriculum & Medium & 0.4442 \\
task\_a\_hybrid\_dynamic\_fusion & Medium & 0.2423 \\
task\_a\_multistage\_hierarchical\_retrieval & Medium & 0.2394 \\
tier1\_cross\_attention\_query\_document\_fixed & Medium & 0.2200 \\
best\_paper\_hierarchical\_routing\_fixed & Medium & 0.2166 \\
tier1\_llm\_distillation & Low & 0.1986 \\
task\_a\_cross\_encoder\_conversation\_context & Low & 0.1949 \\
tier1\_explainable\_retrieval & High & 0.1796 \\
tier1\_causal\_inference & High & 0.1796 \\
tier1\_differentiable\_retrieval & High & 0.1796 \\
tier1\_uncertainty\_aware & High & 0.1796 \\
tier1\_retrieval\_as\_generation & High & 0.1796 \\
\bottomrule
\end{tabular}
\caption{Completed novel experiments with nDCG@10 (retrieval-only).}
\label{tab:completed}
\end{table}

\subsection{Pending Novel Experiments}
Table~\ref{tab:pending} lists implemented novel experiments that are pending execution, along with expected ranges based on planning documents. These are not measured results.
\begin{table}[!t]
\centering
\small
\begin{tabular}{lcc}
\toprule
\textbf{Experiment} & \textbf{Expected nDCG@10} & \textbf{Status} \\
\midrule
CATR & 0.55--0.65 & Pending \\
UQ-Ret & 0.50--0.60 & Pending \\
Diff-Ret & 0.58--0.68 & Pending \\
Meta-Ret & 0.48--0.58 & Pending \\
Causal-Ret & 0.50--0.60 & Pending \\
GEP-Ret & 0.55--0.65 & Pending \\
X-Ret & 0.52--0.62 & Pending \\
CLCF-Ret & 0.53--0.63 & Pending \\
\bottomrule
\end{tabular}
\caption{Pending novel experiments with expected nDCG@10 ranges.}
\label{tab:pending}
\end{table}

\section{Discussion}
The strongest completed novel results come from curriculum learning and hybrid fusion, while high-novelty methods (causal, differentiable, uncertainty-aware) currently show modest nDCG@10. This suggests that novelty must be paired with stronger base models or more robust training to realize its potential. The pending experiments directly target this by combining novel modeling with improved optimization.

\section{Limitations}
This report aggregates available retrieval logs and summaries. Some experiments are pending or stopped, and expected ranges are not measured outcomes.

\section{Conclusion}
We provide a transparent, retrieval-only report of novel ICL-oriented ideas with empirical outcomes, and a readable naming convention for experiments. This structure supports clear review, reproducibility, and a focused path to high-impact results.

\end{document}
